\documentclass{article}
\usepackage{CJKutf8}

\begin{document}

\begin{CJK}{UTF8}{gbsn}
\section*{宋体（简体中文）}
这是一段用于测试简体中文段落排版与自动换行的文本。中文排版要求标点符号不能出现在行首或行尾特定位置，并且汉字之间应当能够正常断行。
\end{CJK}

\bigskip

\begin{CJK}{UTF8}{gkai}
\section*{楷体（简体中文）}
楷体字形流畅，常用于正文中的引用、强调或者书法风格展示。测试楷体中文字符集与字形嵌入。
\end{CJK}

\bigskip

\begin{CJK}{UTF8}{bsmi}
\section*{明體（繁體中文）}
這是一段用於測試繁體中文（傳統漢字）排版與字體嵌入的測試文本。驗證繁體字形如「國」、「學」、「書」、「體」之正確呈現。
\end{CJK}

\bigskip

\begin{CJK}{UTF8}{bkai}
\section*{楷體（繁體中文）}
繁體楷書字體展示，筆意連貫，端莊秀麗。測試繁體楷書在 CJKutf8 下之完整度與嵌入。
\end{CJK}

\end{document}
